\section{Nastanak i izbor AES-a}
\label{sec:nastanak}
% TODO: Razraditi poglavlje tek nakon provere izvora; ne unositi neproverene činjenice.
Ovo poglavlje čini istorijsko jezgro rada. Dokumentacija izbora omogućiće
rekonstrukciju tadašnjih zahteva i obrazloženja izbora, nezavisno od kasnijeg
uspeha standarda~\cite{nist-aes-report-2001,rijndael-proposal-1999}.

\subsection{Potreba za naslednikom DES-a i zahtevi konkursa}
% TODO: Istražiti NIST poziv, rokove, bezbednosne zahteve, javnost procesa i licencne uslove.

\subsection{Kandidati i finalisti}
% TODO: Proveriti spisak kandidata i faze; obraditi MARS, RC6, Rijndael, Serpent i Twofish bez neosnovanog rangiranja.

\subsection{Autori i projektne osobine Rijndaela}
% TODO: Predstaviti Joan Daemen i Vincent Rijmen; koristiti originalni predlog, uz tačnu verziju dokumenta.

\subsection{Kriterijumi i obrazloženje izbora}
% TODO: Razmotriti bezbednost, softverske i hardverske performanse, memoriju, fleksibilnost i različite platforme; izbegavati retrospektivno pripisivanje AES-NI tadašnjim očekivanjima.

\subsection{Izbor 2000. i standardizacija 2001. godine}
% TODO: Dokumentovati tačne datume i FIPS 197; odvojiti uredničko ažuriranje 2023. od tehničke promene algoritma.

\subsection{Rijndael i standardizovani AES}
% TODO: Eksplicitno navesti AES blok od 128 bita i ključeve 128/192/256 bita; razlikovati širu familiju Rijndael sa drugim dimenzijama bloka.
